\section{Current Evidence Base}

\subsection{Fishery Commons as the precursor signal-design study}

The existing Fishery Commons result remains important, but it plays a narrower role in this manuscript than in the earlier two-study paper. Its function is to show that the benchmark can reliably rank central governance signals under repeated adversarial strategy injection. In the matched baseline, monitoring with sanctions reduces held-out collapse relative to no governance under both mutation-based injection and live JSON-based LLM injection. In the broader non-LLM top-down sweep, adaptive quotas emerge as the strongest overall central signal in the medium tier, while temporary closures remain a credible fallback when ecological protection is prioritized over welfare retention. A medium-tier learned-policy validation reaches the same qualitative ordering. Taken together, these results establish that the benchmark can distinguish stronger from weaker central interventions under repeated strategic turnover.

\subsection{Harvest Commons as the current architecture study}

The Harvest Commons result is the more direct architectural finding. The environment extends the commons problem from one shared global stock to local patches with local interaction, local communication, and optional local side-payments. In the current frozen evidence base, the decision-critical comparison is deliberately narrow: top-down-only governance versus hybrid governance under the stronger \texttt{search\_mutation} attacker. Earlier project stages also included \texttt{none} and \texttt{bottom\_up\_only} conditions, but they are not yet the core of the main evidence chain.

Under the current Stage C high-power search-over-mutations matrix, hybrid governance ranks first in all eight decision-critical cells under the manuscript's failure-then-patch-health-then-welfare rule. The clearest consistent effect is ecological rather than catastrophic: garden failure is saturated in the Stage C rows, but hybrid improves patch health and usually reduces neighborhood overharvest relative to top-down-only control. The welfare effect is negative on average but moderate rather than catastrophic. A narrower PPO self-play validation in the medium tier points in the same direction, again primarily through patch-health, local-aggression, and overharvest margins. The current interpretation is therefore that hybrid governance provides better local ecological control under stronger strategic pressure, but not as a free win.

\subsection{What remains incomplete}

Three gaps remain before this architecture result fully supports the broader institutional-alignment framing.

First, the main Study 2 comparison is still too narrow. The environment supports \texttt{none}, \texttt{bottom\_up\_only}, \texttt{top\_down\_only}, and \texttt{hybrid}, but the current high-power evidence only stress-tests the top-down-only versus hybrid pair. That is enough to show that architecture matters, but not enough to say what local-only governance contributes or whether hybrid wins because it combines two useful mechanisms rather than because top-down dominates everything.

Second, the current welfare trade-off is too aggregated. The benchmark already shows that robust governance can impose welfare costs, but it does not yet identify who bears those costs. A stronger institutional reading requires that the analysis distinguish whether governance primarily penalizes highly exploitative agents, broadly taxes everyone, or redistributes welfare in more complex ways.

Third, attacker strength is present in the codebase but not yet fully organized in the story. The existing pipeline already contains a ladder from random injection through mutation and heuristic attackers to search over mutations, with a separate diagnostic LLM path. However, the architecture result is not yet presented as a capability ladder that asks when each governance arrangement breaks.
